\documentclass[11pt]{article}
\usepackage{amsmath,amssymb,amsthm,fullpage}
\newtheorem{lemma}{Lemma}
\title{Curved Poles and Recursive Cauchy Cancellation}
\author{}
\date{November 24, 2026}
\begin{document}
\maketitle

\paragraph{Verification status.}
The finite four-pole construction below has rational witnesses for every
claimed strict sign table. The recursive cancellation identities have
complete written proofs. Neither proves a superconstant saving in
HeadComplexity. In particular, simultaneous vanishing is not a proof of
simultaneous sign interpolation. No model definition is changed and no
claim in this note is Lean checked.

\section{An exact escape from the fixed-pole obstruction}
Let $\epsilon=10^{-5}$, $\eta=10^{-3}$, and
\[
 p_h=\epsilon(4+a_h)+16\epsilon^2\mathbf1_{h=1},
 \qquad (a_0,a_1,a_2,a_3)=(-3,-1,1,3).
\]
Take eight times
\[
 t_i=1+4\epsilon+\eta u_i,
 \qquad (u_0,\ldots,u_7)=(-3,-1,1,3,-2,2,0,0)
\]
and masks
\[
 S_0=\cdots=S_3=\{0,1,2,3\},\quad
 S_4=S_5=\{1,3\},\quad S_6=\{2\},\quad S_7=\{3\}.
\]
For each exponent $\nu\in\{1,2\}$ and each desired eight-bit
strict sign table, there are a permutation $\pi$ of the four poles
and four rational residues $r_h$ realizing that table with
\[
 f_i=\sum_{h:\pi(h)\in S_i}\frac{r_h}{(t_i-p_h)^\nu}.
\]
The companion program \texttt{curved-pole-tree-check.py} constructs
and checks all $256$ witnesses for each exponent. Its JSON output
records each permutation, each rational residue, and every exact
score. It checks the strict signs directly in the original matrix;
no numerical infeasibility assertion is involved.

This construction changes a pole by order $\epsilon^2$, so it is
outside the fixed affine pole family in
\texttt{butterfly-tree-cauchy-audit.tex}. It disproves any extension
of that obstruction to arbitrary varying poles. It does not show a
superconstant gain: the eight signs are no more than the eight real
pole and residue parameters.

\section{The reciprocal coordinate for simultaneous cancellation}
The natural pole coordinate near an exterior query time $T$ is
\[
 q_h=\frac1{T-p_h}.
\]
For positive distinct $q_h$, choose $p_h=T-1/q_h$. Taking all $q_h$
slightly above $1/T$ makes every $p_h$ positive and as small as desired.
If $Q(q)=\prod_{h=1}^M(q-q_h)$ and
\[
 w_h=\frac1{Q'(q_h)},\qquad r_h=\frac{w_h}{q_h},
\]
then the exterior score satisfies
\[
 \sum_h\frac{r_h}{T+s-p_h}
 =\sum_h\frac{w_h}{1+s q_h}
 =\frac{(-s)^{M-1}}{\prod_h(1+s q_h)}.
 \tag{1}\label{eq:barycentric}
\]
The last identity follows from the partial fraction identity
$1/Q(u)=\sum_h w_h/(u-q_h)$, substituted at $u=-1/s$.
It extends to $s=0$ by continuity. Thus the score has exactly order
$M-1$ at $s=0$ when $M>1$.

\begin{lemma}[Restriction to constant complementary factors]
Let a finite simple root set be partitioned into two children,
whose monic polynomials are $A(q)$ and $B(q)$. Suppose
$B(q)-A(q)=c\ne0$. At every root $q_h$ of $A$,
\[
 \frac1{(AB)'(q_h)}=\frac1{c A'(q_h)},
\]
and at every root of $B$ the corresponding identity has factor
$-1/c$ and derivative $B'$. Consequently the full-set barycentric
weights restricted to either child are proportional to that child's
barycentric weights.
\end{lemma}
\begin{proof}
At a root of $A$, the derivative of $AB$ is $A'B=cA'$.
At a root of $B$, one has $A=-c$ and $(AB)'=-cB'$.
\end{proof}

Apply this lemma recursively to a binary tree of polynomial root
sets whose two monic child polynomials differ by a nonzero constant.
The global barycentric weights restricted to every node are a
nonzero constant times that node's own barycentric weights. Equation
\eqref{eq:barycentric} then proves simultaneous vanishing to order
$M-1$ for each node of size $M$.

Such real simple root sets exist at every depth. The roots of a
Chebyshev polynomial $T_{2^d}(x)-c$, with $-1<c<1$, split using
\[
 T_{2m}(x)=2T_m(x)^2-1.
\]
The children are roots of
$T_m(x)\mp\sqrt{(c+1)/2}$; after the common leading coefficient
is removed, their monic polynomials differ by a nonzero constant.
Every resulting parameter remains strictly between $-1$ and $1$,
and all roots are simple and real. An affine change of $x$ places
the reciprocal coordinates in any desired small positive interval.

One rational depth-three example uses the nodes
$\{\pm1,\pm4,\pm7,\pm8\}$. Its root polynomial is
\[
 (x^4-65x^2+424)^2-360^2.
\]
The children are $\{\pm1,\pm8\}$ and $\{\pm4,\pm7\}$,
which in turn split into their symmetric pairs. The identity
$1^2+8^2=4^2+7^2=65$ makes the two degree-four polynomials
differ by a constant.

\section{The precise missing interpolation statement}
For $H=2^d$ leaves, give the root $H$ full-mask queries and every
right-child node of size $M$ exactly $M$ queries. There are
\[
 H+\frac{Hd}{2}
\]
queries. The hoped-for statement would realize every strict sign
table by assigning pole and residue pairs to the leaves, while
keeping all query times fixed and close to $T$ and every pole small.
The cancellation lemma supplies only a singular base configuration.
It does not provide the required signs after perturbation.

For a fixed root set, the exact compatible binary partitions are
severely constrained. Indeed if monic children satisfy $B-A=c$,
then $AB=((A+B)/2)^2-c^2/4$. Comparing the coefficients in the
upper half of this square determines the monic polynomial $(A+B)/2$
uniquely from the parent polynomial; its constant term is then also
determined by the next coefficient whenever the child degree is
positive. Equivalently, two monic polynomials of the same positive
degree cannot have squares differing by a nonzero constant, since
their sum has positive degree and their difference would have to
multiply it to a constant. Hence the compatible child partition is
unique up to swapping children. Recursively, such swaps give only
$2^{H-1}$ assignments for a fixed compatible root set. Counting
these assignments as $H!$ independent choices would be invalid.

The remaining issue is whether label-dependent pole deformations,
or partitions with fewer simultaneous cancellations, supply enough
additional freedom for all $H+Hd/2$ strict signs. Neither the finite
four-pole check nor the barycentric identities answer this issue.

\end{document}
